\documentclass[11pt]{article}

\usepackage[T1]{fontenc}
\usepackage[utf8]{inputenc}
\usepackage{lmodern}
\usepackage[margin=1in]{geometry}
\usepackage{amsmath,amssymb,amsthm,mathtools,bm}
\usepackage{longtable,booktabs,array,tabularx}
\usepackage[shortlabels]{enumitem}
\usepackage{xcolor}
\usepackage{hyperref}
\usepackage{xurl}
\usepackage{microtype}
\usepackage{seqsplit}

\hypersetup{colorlinks=true,linkcolor=blue,citecolor=blue,urlcolor=blue}
\hypersetup{pdftitle={The Fine-Structure Constant as an OPH Pixel Fixed Point},pdfauthor={Bernhard Mueller; Alexander Osika; Daniel Ari Friedman}}
\setlength{\parindent}{0pt}
\setlength{\parskip}{0.55em}
\setlength{\emergencystretch}{3em}
\renewcommand{\arraystretch}{1.12}
\newcolumntype{L}[1]{>{\raggedright\arraybackslash}p{#1}}
\newcolumntype{R}[1]{>{\raggedleft\arraybackslash}p{#1}}

\newtheorem{theorem}{Theorem}
\newtheorem{definition}{Definition}
\newtheorem{proposition}{Proposition}
\theoremstyle{remark}
\newtheorem{remark}{Remark}

\newcommand{\OPH}{\mathrm{OPH}}
\newcommand{\U}{\mathrm{U}}
\newcommand{\SU}{\mathrm{SU}}
\newcommand{\A}{\mathcal A}
\newcommand{\R}{\mathbb R}
\newcommand{\Z}{\mathbb Z}
\newcommand{\dd}{\,\mathrm d}
\newcommand{\e}{\mathrm e}
\newcommand{\alphainv}{\alpha^{-1}}
\newcommand{\ATh}{A_{\mathrm{Th}}}
\newcommand{\AZ}{A_Z}
\newcommand{\Pstar}{P_\star}
\newcommand{\Ppub}{P_{\mathrm C}}
\newcommand{\Rq}{R_Q}
\newcommand{\Dcalc}{\Delta_{\mathrm{calc}}}
\newcommand{\Dcalib}{\Delta_{\mathrm{H,cal}}}
\newcommand{\Dlep}{\Delta_{\mathrm{lep}}}
\newcommand{\Dhad}{\Delta_{\mathrm{had}}}
\newcommand{\DEW}{\Delta_{\mathrm{EW}}}
\newcommand{\Dreq}{\Delta_{\mathrm{req}}}
\newcommand{\Dhadreq}{\Delta_{\mathrm{H,req}}}
\newcommand{\Dmiss}{\Delta_{\mathrm{miss}}}
\newcommand{\Nc}{N_c}
\newcommand{\Ng}{N_g}
\newcommand{\decnum}[1]{\vcenter{\hbox{\parbox{0.68\linewidth}{\small\ttfamily\seqsplit{#1}}}}}

\definecolor{ophblue}{HTML}{1B4E7A}
\definecolor{ophgreen}{HTML}{1F7A5A}
\definecolor{ophgold}{HTML}{A66A00}
\definecolor{ophred}{HTML}{8A2D2D}

\title{\textbf{The Fine-Structure Constant as an OPH Pixel Fixed Point}}
\author{Bernhard M\"uller \and Alexander Osika \and Daniel Ari Friedman}
\date{\OPHPaperReleaseDate}
\newcommand{\OPHPaperReleaseID}{r2056}
\newcommand{\OPHPaperReleaseDate}{October 31, 2026}
\newcommand{\OPHPaperReleaseBanner}{%
\begin{center}
\small\textbf{Paper release:} \texttt{\OPHPaperReleaseID}\quad\textbf{Released:} \OPHPaperReleaseDate
\end{center}
}
\newcommand{\PaperReleaseID}{\OPHPaperReleaseID}
\newcommand{\PaperReleaseDate}{\OPHPaperReleaseDate}
\newcommand{\PaperReleaseBanner}{%
\begin{center}
\small\textbf{Paper release:} \texttt{\PaperReleaseID}\quad\textbf{Released:} \PaperReleaseDate
\end{center}
}


\begin{document}
\maketitle
\OPHPaperReleaseBanner

\begin{center}
\small\textbf{Author affiliations:} Bernhard Mueller, Pragma Research Inc.; Alexander Osika, EtherWorks; Daniel Ari Friedman, Active Inference Institute and FractiAI.
\end{center}

\begin{abstract}
This paper studies a conditional pixel fixed point for the fine-structure
constant. The same-cell equation \(P=\varphi+\sqrt\pi\,\alpha(0)\) relates a
screen-area coordinate to an electromagnetic readout, and a declared source
map returns the electromagnetic readout from the same coordinate through
declared closure coordinates and running conventions. Two declared trial
maps have interval-certified unique fixed points, with inverse couplings
\(136.994835177413\ldots\) and \(137.035660136947\ldots\), about 300 and
2.5 parts per million from the CODATA 2022 value. The first transports the
electroweak anchor to the Thomson endpoint with a declared quark-screening
ansatz; the second adds the finite-screen unified gauge width inside the
loop. Replacing the ansatz by the measured hadronic spectral payload gives
\(136.3827548\) on the interval \([136.3670481,136.3984652]\), which excludes
the measured value; reaching it would need a hadronic shift of \(0.032244\)
instead of the measured \(0.027609\pm0.000112\), or a same-scheme source-anchor
bridge of about \(0.635\) inverse-alpha units. The calculation is a
conditional diagnostic, not a same-process derivation from the executable
model: its inputs include a supplied charged-fermion spectrum, a declared
gauge-running convention and measured hadronic spectral data.

An independent entangled-pair construction due to A.~D.~Wissner-Gross gives
the exact scalar \((79+48\sqrt2+42\sqrt3+3\sqrt6)/31104\), whose inverse is
\(137.035999189566\ldots\). We reconstruct its finite exchange/projector
arithmetic and complete measurement controls; if its conjectured Thomson
identification and the same-cell reading are granted, it supplies another
conditional area coordinate. The constructions supply no source-selected
photon normalization, hadronic threshold transport, absolute particle masses
or physical same-cell bridge, and every numerical endpoint comparison is a
retrospective diagnostic.
\end{abstract}

\noindent\textbf{Keywords:} fine-structure constant, Observer-Patch Holography, pixel fixed
point, electromagnetic coupling, Standard Model constants, hadronic vacuum polarization

% Shared statement of the OPH tower and its executable hypothesis.
% Included by every active paper. No theorem premise or physical status is
% promoted by this panel; each paper states its own layer and hypothesis.
\paragraph{One theory and its executable hypothesis.}
OPH proposes closure: one complete history is both a timeless
self-referential structure and its own internal computation, and observers
inside it recover and realize its specification. The hypothesis that makes
this proposal concrete and testable is executable.

\begin{quote}\itshape
A simple, fully specified local repair and memory law on the proposed
carrier architecture admits continuing observers whose compatible records
and interventions have the observed \(3+1\)-dimensional quantum-field
description and a common classical-gravity limit. The preferred matter
mechanism is localized topological defects of native overlap transport,
with rest energy derived from their excess maintenance dynamics in the same
physical clock and energy dictionary.
\end{quote}

One source process must meet every part of this statement: records, clocks,
matter, energy and responsive geometry come from the same law, state and
refinement. Coherence, protected codes, radial directions, fields, clocks
and couplings may be specified as part of its design; their implementation
and observable consequences must be established on that process. A failed
specified law rejects that law and its tested interpretation; it rejects the
wider proposal only when the rejected choice is proved necessary.

\paragraph{The tower.}
The papers form one tower over the executable model. Each layer states a
hypothesis and proves finite theorems under displayed premises. A conclusion
passes to the next layer only after that premise is proved on the same
model process; until then it is a conditional implication.
\begin{description}\setlength{\itemsep}{2pt}
\item[0. Closure and agreement.] One history reads itself into existence, and
  agreement on shared records defines the public world. This is the
  interpretation and acceptance criterion; no theorem selects a unique
  universe.
\item[1. Repair mathematics.] Local agreement repair reaches global
  consistency, and nonzero cycle holonomy is the residue that local repair
  cannot remove. Fixed-point, normal-form, frustration and finite
  event-algebra theorems hold on their stated domains.
\item[2. Carrier and repair law.] Federated twelve-port icosahedral carriers
  with the native C/E seam repair are the substrate. The seam rule follows
  from four stated requirements, the quantum record laws are classified, and
  the uniform native repair selects a rank-three slow-response band.
\item[3. Executable observer.] Observers are continuing native processes that
  prepare, read, retain and act on records. Finite complete-output
  executions, overlapping-observer correspondence, event clocks and
  checkpoint continuation are verified; one continuing observer with every
  service and resource accounted is the open join.
\item[4. Clocks and fields.] Modular flow of observer algebras supplies
  candidate local clocks, and native repair supplies a positive continuum
  generator of spectral dimension two on the screen. A declared radial
  extension gives a conditional smooth \(1+3\) scalar description. Complete
  finite source factors join spatial CAR dynamics, particle creation,
  repair and a common calibrated clock. Physical validation of that clock
  and energy normalization, and an interacting chiral \(3+1\) limit, remain
  open.
\item[5. Gauge structure and matter.] Gauge structure is the transport of
  overlap data, and stable massive matter is a localized defect whose
  boundary transport carries a holonomy that admitted repairs cannot erase.
  The compact Lie type \(\mathfrak u(1)\oplus\mathfrak{su}(2)\oplus
  \mathfrak{su}(3)\) follows under complete-response premises, and fixed
  nonzero holonomy forces positive residual mismatch. A specified coherent
  ordering memory protects native exchange-history sectors, and a colored
  extension implements their fermionic exchange. A localized physical
  defect, its maintenance-to-mass law and a predicted measured spectrum
  remain unestablished; the earlier double-load candidate is transient.
\item[6. Gravity.] Einstein dynamics is the coarse response of the same
  process. The null-balance and tensor-completion implication is proved
  under its premises; realizing those premises on one refinement tower is
  open, and the static radial product fails null balance.
\item[7. Cosmology and tests.] The same process must fix the vacuum, dark
  sector, primordial response and black-hole records, and pass independent
  prospective tests. The cosmology papers and the fine-structure, Koide,
  de Sitter and Yang--Mills supplements carry conditional branches,
  diagnostics and retrospective comparisons at this layer.
\end{description}

\paragraph{Completion and stopping.}
The aim is a consistent and complete theory: every layer realized on one
process, or a proved obstruction that removes OPH or a stated branch of it.
Results that no join needs leave the active chain with their provenance. A
universal obstruction over a frozen admissible class, an unavoidable
excluded prediction, or a contradiction among indispensable premises stops
that class. A failed candidate, an unconstructed map or a resource stop has a
narrower meaning. The ordered joins, retained foundations and stopping
rules are recorded with the executable model in the public
repository.\footnote{\url{https://github.com/FloatingPragma/observer-patch-holography/blob/r2056/docs/COMPLETION_PROGRAM.md}}

\paragraph{Fixed physical boundaries.}
Repair counts, host runtime and entropy production are distinct from
energy; the maintenance-to-energy map must be derived with one common
signal speed and no per-species conversion. The modular parameter is
dimensionless until a specified and implemented clock calibration relates
it to duration. Stationary repair or
record entropy is distinct from the fixed-volume generalized-entropy
variation that the Einstein construction consumes. A massless photon makes
the infinite-volume theory gapless, so sector gaps and thresholds name their
state and limit; confinement precludes free massless gluons as observable
particles; a quantum graviton lies outside the classical Einstein theorem.


\section*{This paper's hypothesis and place in the tower}

This paper carries a conditional branch at layer seven. Its hypothesis is a
pixel fixed point: if a typed bridge proves that the detuning
\((P-\varphi)/\sqrt\pi\) of one screen cell from the golden-ratio balance
and the Thomson-limit coupling \(1/\ATh(P)\) returned by a declared source
map at the same \(P\) read one cell quantity, then closure requires
\(P=\varphi+\sqrt\pi/\ATh(P)\), and \(\alpha(0)\) is the coupling at the
fixed point. The source map consumes declared closure coordinates and
conventions: the golden-ratio balance point, the Gaussian width
\(\sqrt\pi\), the \(P/4\) representation-entropy budget, the one-loop MSSM
coefficient triple, the transmutation coefficient \(\beta_{\rm EW}=4\), a
supplied phenomenological charged-fermion spectrum and a declared
quark-screening ansatz. The measured hadronic spectral payload enters only
as a replacement test of that ansatz, and the CODATA value only as the
endpoint.

The finite results are exact: each declared map has one interval-certified
fixed point on its maximal analytic domain; coarse register data do not
determine the hadronic residual, and a 24-site Jacobi realization of a
chosen residual is tautological; the conditional pure-source theorem is
Banach's theorem under its listed premises. The calculation is not a
same-process derivation from the executable model. That model supplies no
physical area map for its screen cells, no Lorentzian Maxwell theory, no
charged spectrum and no hadronic spectral measure, and it uses the source
root only as the closure input of one detector.

The open joins are a target-independent rule selecting one map, the typed
same-cell bridge, a native area and clock dictionary for \(P\), charged
masses from native defects, and a source-emitted Ward-current hadronic
spectral measure with its same-scheme remainder. With the measured hadronic
payload in place of the ansatz, the declared transport excludes the
measured value, so the branch needs a same-scheme anchor bridge of about
\(0.635\) inverse-alpha units. A target-independent selected map whose
certified prediction interval, narrower than the direct rubidium
measurement, excludes the measured value would refute the branch.

\section{Introduction}

The fine-structure constant is usually treated as an empirical constant.
Here it is studied as a conditional readout of one screen cell. On the
declared carrier branch the cell belongs to the federated screen built from
twelve-port icosahedral carriers, the structure whose port algebra
witnesses the Standard Model gauge structure in the gauge
paper~\cite{oph-smgr}. The known physics is the low-energy Maxwell coupling,
charged-lepton running and low-energy hadronic spectral transport.

Two declared principles carry the construction. The detuning law places the
pixel ratio off the golden-ratio balance point by the Gaussian-normalized
width of one electromagnetic observation, \(P=\varphi+\sqrt\pi\,\alpha\).
The constant-identification law is a closure premise: a typed bridge must
prove that the fine-structure constant read inside the trial world and the
substrate pixel readout are one quantity. Conditional on that bridge, the
two principles leave one local variable and one self-read equation. Neither
the fixed-point mathematics nor the source certificates construct the bridge
or select between the two certified maps. The measured fine-structure
constant enters only after the source calculation, as an endpoint test.

Jentschura and Nandori survey first-principles attempts to understand
\(\alpha\), including beta-function, symmetry, cutoff and string-inspired
ideas \cite{jentschura-nandori}. Golden-ratio constructions appear in
geometric proposals \cite{heyrovska-narayan} and in semi-empirical work on
electroweak and flavor mixing \cite{ciborowski}. Octonionic and
exceptional-algebra approaches also produce values near \(1/137\)
\cite{singh}, and symbolic regression and quantum-information criteria have
been used to look for structure in Standard Model constants and electroweak
parameters \cite{chekanov-kjellerstrand,liu-low-yin}. The calculation here
belongs to this search for structure. Its physical interpretation requires a
target-independent choice of the map, a typed bridge proving that its input
and output read the same local screen cell, and same-scheme transport to the
laboratory endpoint.

\begin{quote}\small
\textbf{Certified source roots.} The executable source map contracts to the
undressed source/root inverse coupling
\(\alphainv_{\mathrm{root}}=136.994835177413\ldots\), the
interval-certified unique fixed point of the declared map (public interval
certificate \cite{oph-code}; enclosure width \(7.2\times10^{-24}\),
\(L\le0.0724\)), 300.388218 parts per million below the CODATA 2022 value
\cite{ophpredictions}. Its hadronic transport is the declared
quark-screening ansatz \(S_{\mathrm{calc}}\approx0.887\) of Step~14, which
carries no perturbative derivation. The self-consistent map that adds the
finite-screen unified gauge width \(\alpha_U(P)\) inside the loop has the
certified fixed point \(137.035660136946577\ldots\), with relative residual
\(2.5\times10^{-6}\), about \(1.6\times10^{4}\) measurement sigma. The
combination \(A_{\alpha_U}^{\mathrm{fp}}=137.035959513609\ldots\) mixes the
certified source root with \(\alpha_U(\Ppub)\) evaluated at the
CODATA-derived comparison pixel; it is not a fixed point of a single
declared map.

\textbf{CODATA\slash NIST measured endpoint.} CODATA\slash NIST reports the
value \(137.035999177(21)\) for \(\alphainv(0)\). The difference from
\(A_{\alpha_U}^{\mathrm{fp}}\) is CODATA-minus-diagnostic bookkeeping; it is
not the output of a same-scheme hadronic spectral calculation.
\end{quote}

The proximity of both roots to the measured value depends on the screening
ansatz: replacing it by the measured hadronic payload gives
\(136.3827548\), whose certified interval excludes the measured value
(Section~\ref{sec:pdg-step}). The residual is therefore not explained by
missing hadronic transport, and the comparison does not identify the
physical source of the gap.

\paragraph*{Selection exposure and numerical support.}
Two declared maps have interval-certified fixed points. A separate
asymptotic compressed calculation prints \(136.9940207\), but is an
approximate candidate: its printed grain and inverse coupling disagree in
the outer identity by about \(5.25\times10^{-6}\) relatively. It is not a
third certified root. A separate retrospective constant-substitution grid
contains 48 labels, 42 distinct numerical pairs and 41 alternatives to the
canonical pair; none of those alternatives reaches the declared relative
distance \(2.5\times10^{-6}\). That threshold was chosen after the canonical
distance was known. These counts describe numerical coverage, not a
probability of accidental agreement. No exhaustive map-form menu or
sampling measure is specified. Even the same two certified outputs,
with one inside that distance threshold, give hit probability \(1/2\) under
equal weights and \(3/4\) under weights \((1/4,3/4)\). Those are mathematical
counterexamples to identifying significance from output values and counts;
neither weighting is inferred for this calculation. The measured endpoint
therefore assigns no calibrated chance probability to the choice of map.

\paragraph*{Provenance of the three numbers.}
The chain
\[
136.994835177413\ldots
\longrightarrow
137.035660136946577\ldots
\longrightarrow
137.035999177(21)
\]
runs from the source pixel through the gauge-width dressed pixel to the
measured low-energy value. The first two values are the interval-certified
fixed points of the declared source map and of the declared self-consistent
gauge-width map \cite{oph-code}; the third is the measured endpoint. The
root value \(\alpha_{\mathrm{root}}^{-1}\) is returned by the incomplete
source map before the low-energy charged vacuum is attached. The gauge width
\(\alpha_U\) is the main high-scale gauge-sector width carried by the pixel:
the weak and color representation content must fit the pixel budget
\(P/4\), and the electroweak heat-kernel closure supplies the executable
certificate. It supplies almost all of the inverse-alpha gap from the naked
source value.

The factor \(C_{24,Q}\) of Step~18 is a back-solved calibration: the small
endpoint dressing required after the gauge-width contribution is added,
solved from the measured endpoint rather than computed from source data.
Its excess \(\alpha_U(\Ppub)(C_{24,Q}-1)\) is a comparison residual with the
numerical size required of the residual same-scheme low-energy hadronic
endpoint transport; it does not compute that transport. A zero-momentum
laboratory photon sees the Ward-projected \(\mathrm U(1)_Q\) current after
charged-lepton vacuum polarization, confined-quark/hadron spectral
transport, and finite endpoint matching, so the observer measures the fully
dressed low-energy coupling, while \(A_{\alpha_U}^{\mathrm{fp}}\) is a
mixed-provenance no-hadron diagnostic before that dressing.

The observer account is the flagship synthesis \cite{oph-synthesis}, the
source map and fixed-point branch are recorded in the Standard Model gauge
paper \cite{oph-smgr}, the charged-spectrum and particle-structure
continuation in the particle paper \cite{oph-particle}, and the
positive-chamber balance identity used by the charged-lepton continuation in
the Koide paper \cite{oph-koide}.

\begin{proposition}[Maxwell normalization of the endpoint branch]
On the declared electromagnetic branch used by the fixed-point equation, the Ward-projected
\(\U(1)_Q\) channel carries
\[
F_Q=dA_Q,\qquad dF_Q=0,\qquad d{*}F_Q=g_Q^2(q^2;P){*}J_Q,
\]
with
\[
g_Q^2(q^2;P)=4\pi\alpha_{\mathrm{em}}(q^2;P),
\qquad
\ATh(P)=\alpha_{\mathrm{em}}^{-1}(0;P).
\]
Thus the inside reading in the pixel equation is the Thomson-limit Maxwell coupling of the same
electromagnetic current.
\end{proposition}

\begin{proof}
On the declared unbroken Maxwell action of the gauge paper, the electromagnetic branch is the
abelian factor \(\U(1)_Q\). There the compact-gauge curvature reduces to \(F_Q=dA_Q\), and the
quadratic field action gives \(d{*}F_Q=g_Q^2{*}J_Q\). That action is a branch premise: the
executable observer supplies finite Maxwell-shaped seam actions, and a local source-produced
physical Maxwell theory on Lorentzian spacetime is not constructed \cite{ophpostdiction}. The particle-paper transport theorem then identifies
the same \(g_Q(q^2;P)\) with the Ward-projected running electromagnetic coupling. Taking
\(q^2\to0\) gives the endpoint quantity \(\ATh(P)\) used in the pixel fixed-point equation.
\end{proof}

The hadronic step uses the standard dispersion logic behind hadronic vacuum-polarization work. The relevant comparison data are electromagnetic spectral functions measured through \(e^+e^-\to\mathrm{hadrons}\), as used in data-driven running-\(\alpha\) and \(g-2\) analyses \cite{dhmz,knt19,pdg-xsect}. In this paper, the decimal correction shown above is reported as the difference between the CODATA/NIST comparison value and the OPH pure source calculation. The source data in this paper emit no direct production calculation of the hadronic spectral function; a pure source endpoint theorem consumes that calculation as a premise.

\clearpage
\section{Symbols}

\begin{longtable}{@{}L{0.20\textwidth}L{0.32\textwidth}L{0.38\textwidth}@{}}
\toprule
Symbol & Definition & Informal meaning \\
\midrule
\endhead
\(\A(P)\) & Patch algebra assigned to a screen patch \(P\) & The local observables available to one finite observer patch. \\
\(\omega_P\) & State on \(\A(P)\) & The local physical description carried by that patch. \\
\(\A(P\cap Q)\) & Shared overlap algebra & The observables two neighboring patches can compare. \\
\(P\) & \(a_{\mathrm{cell}}/\ell_\star^2\) & The dimensionless area of one screen cell in scale-readout units. \\
\(a_{\mathrm{cell}}\) & Declared screen-cell area & The area assigned to the local pixel by the declared screen dictionary. \\
\(\ell_\star\) & OPH scale-certificate length & The length emitted by the selected no-\(G\) scale certificate. \\
\(\ell_P\) & Planck length & The SI display of \(\ell_\star\) after \(G_{\mathrm{SI}}=c^3\ell_\star^2/\hbar\). \\
\(\varphi\) & \((1+\sqrt5)/2\) & The self-similar entropy-balance point. \\
\(\sqrt\pi\) & Boundary Gaussian normalization width & The conversion factor from pixel displacement to observation width. \\
\(\alpha_{\mathrm{ext}}(P)\) & \((P-\varphi)/\sqrt\pi\) & The declared outer detuning coordinate. \\
\(\ATh(P)\) & \(\alpha_{\mathrm{em}}^{-1}(0;P)\) & The inverse electromagnetic endpoint returned by the declared trial map. \\
\(\alpha_{\mathrm{in}}(P)\) & \(1/\ATh(P)\) & The declared inner coordinate; its same-cell interpretation requires a typed bridge. \\
\(E_P\) & Planck energy & The energy unit associated with the selected scale certificate. The numerical equations use Planck units unless GeV labels are attached. \\
\(M_U(P)\) & \(E_P\e^{-2\pi}P^{1/6}\) & The OPH unification-scale readout returned at the declared trial coordinate. \\
\(E_{\mathrm{cell}}(P)\) & \(E_P/\sqrt P\) & The local cell energy readout. \\
\(\alpha_U(P)\) & Unified coupling solved from the heat-kernel closure equation & The coupling at \(M_U(P)\) that makes the gauge representation entropy match the pixel. \\
\(\alpha_i(\mu;P)\) & Gauge coupling \(i=1,2,3\) at scale \(\mu\) & The hypercharge, weak, and color couplings emitted by the same trial \(P\). \\
\(b_i\) & \((33/5,1,-3)\) & The one-loop MSSM beta coefficients used as the high-scale running convention here; a declared structural selection. \\
\(m_Z(P)\) & Self-consistent electroweak anchor scale & The \(Z\)-scale generated by \(v(P)\), \(\alpha_1(P)\), and \(\alpha_2(P)\). \\
\(v(P)\) & \(E_{\mathrm{cell}}(P)\exp[-2\pi/(\beta_{\mathrm{EW}}\alpha_U(P))]\) & The electroweak transmutation scale. \\
\(\beta_{\mathrm{EW}}\) & \(\Nc+1=4\) for \(\Nc=3\) & The coefficient that controls the exponential drop to the electroweak scale. \\
\(\AZ(P)\) & \(\alpha_{\mathrm{em}}^{-1}(m_Z^2;P)\) & The electroweak-scale electromagnetic anchor. \\
\(\Dlep(P)\) & Lepton transport contribution to inverse alpha & The exact one-loop kernel evaluated on the frozen phenomenological \(e,\mu,\tau\) spectrum. \\
\(\Dhad(P)\) & Hadronic transport contribution & The low-energy QCD spectral contribution needed for the pure source calculation. \\
\(\DEW(P)\) & Electroweak/scheme endpoint remainder & Same-scheme finite remainder needed to connect the source anchor to the endpoint. \\
\(\Dcalc(P)\) & Calculated charged-fermion transport contribution & The exact kernel on the supplied lepton continuation plus the screened quark continuation; not a pure source prediction. \\
\(\Dcalib^{\mathrm{fp}}\) & Root-to-CODATA bookkeeping difference & The inverse-alpha amount between the naked root value and the CODATA/NIST central endpoint; it includes the comparison-pixel \(\alpha_U(\Ppub)\) contribution and is not the post-\(\alpha_U\) hadronic gap. \\
\(\Dhadreq^{\mathrm{fp}}\) & Required hadronic endpoint correction after \(A_{\alpha_U}^{\mathrm{fp}}\) & The inverse-alpha amount by which the displayed OPH output misses the CODATA/NIST central endpoint. \\
\(\Rq(P)\) & \(\ATh(P)-[\AZ(P)+\Dcalc(P)]\) & The residual same-scheme hadronic endpoint contribution. \\
\(A_{\alpha_U}^{\mathrm{fp}}\) & \(\alpha_{\mathrm{root}}^{-1}+\alpha_U(\Ppub)\) & Mixed-provenance no-hadron comparison diagnostic; the root is source-side and \(\Ppub\) is CODATA-derived. \\
\(J_{24,Q}(P)\), \(\omega_Q(P)\), \(\Xi_Q(P)\) & Source-side hadronic spectral objects required for a final endpoint theorem & The Jacobi/Stieltjes spectral payload, normalization, and finite same-scheme remainder that must be emitted without looking at the Thomson target. \\
\(\Ppub\) & CODATA/NIST comparison pixel & The pixel obtained from the CODATA/NIST Thomson endpoint under the OPH outer equation. \\
\bottomrule
\end{longtable}

\section{Shared Collar Convention for \(\chi_\nu\)}

The \(\chi_\nu\) protected-reserve theorem uses the alpha-calibrated comparison
pixel
\[
P_\chi=P_{\mathrm C}=1.630968209403959\ldots .
\]
Its scalar reserve density is
\[
\epsilon_{\rm res}=\frac{P_\chi}{24}
\]
only after the same-collar shared-edge budget and one-class \(\mathbb Z_6\)
trace receipts pass in the screen microphysics paper and in the observers library paper.  This
does not make \(P/4\) a primitive Hilbert-space dimension.  In this paper,
\(P/4\) is the local screen-cell entropy budget used by the electromagnetic
fixed-point branch; \(P_\chi/24\) is the shared scalar-reserve density used by the
\(\chi_\nu\) collar branch. The \(P/4\) budget is a declared area reading of the cell.
Native regional entropy in the executable observer grows with the number of enclosed ports
rather than with crossing seams, and the physical region map that would carry an area law is a
premise \cite{ophentropy}.

\section{Step 1: The OPH Overlap Rule}

OPH starts from a finite screen covered by observer patches. If \(P_1\) and \(P_2\) overlap, then their induced states must agree on the shared algebra:
\begin{equation}
\omega_{P_1}|_{\A(P_1\cap P_2)}
=
\omega_{P_2}|_{\A(P_1\cap P_2)}.
\label{eq:overlap}
\end{equation}

Informally: no observer sees the whole world. Physics is what survives when neighboring local descriptions can be checked against each other and made consistent on the parts they share.

The quantitative fine-structure branch applies this rule to a single local
screen cell. The cell has one outside description and one inside description.
The outside description is geometric. The inside description is
electromagnetic. A declared typed bridge identifies both as readings of one
cell; the fixed-point theorem is conditional on that premise.

The executable observer construction runs local agreement repair on a twelve-port carrier:
each native repair event moves load across one seam and records its placement coin and the
transferred current \cite{ophmodel}. Its physical identification with our Universe is unestablished.

Why this step is needed: OPH reads a dimensionless constant as a closure coordinate that finite observers locate by measurement. A declared map is then tested by comparing its certified fixed point with the measured value, and that comparison is a diagnostic.

\section{Step 2: The Pixel Variable}

Define the local pixel ratio
\begin{equation}
P:=\frac{a_{\mathrm{cell}}}{\ell_\star^2}.
\label{eq:Pdef}
\end{equation}
Here \(a_{\mathrm{cell}}\) is the screen-cell area and \(\ell_\star^2\) is the area scale emitted
by the separate OPH gravity readback fixed point.

Informally: \(P\) says how large the screen cell is when measured in scale-readout units. It is dimensionless, so it can be compared directly with pure numbers such as \(\varphi\) and \(\sqrt\pi\). After the gravity row emits \(G\), the same scale is displayed as the Planck area in SI units. The pixel equation fixes this dimensionless ratio; it does not determine the SI scale by itself.

Why this step is needed: \(\alpha\) has no units. A unitless electromagnetic number has to be matched to a unitless geometric number. Dividing the cell area by the emitted scale area gives that number without using the pixel equation to derive \(G\).
The executable observer construction supplies no physical area map for its screen cells, so this
area reading of \(P\) is a declared dictionary entry \cite{ophconstruction}.

\section{Step 3: The Golden-Ratio Balance}

The self-similar balance point is
\begin{equation}
\varphi=1+\frac{1}{\varphi},
\qquad
\varphi^2-\varphi-1=0,
\qquad
\varphi=\frac{1+\sqrt5}{2}.
\label{eq:phi}
\end{equation}

Informally: \(\varphi\) is the fixed point of the simplest self-similar split, where the whole-to-large ratio equals the large-to-small ratio. OPH takes it as the declared zero-detuning balance point for the local pixel.

The pixel does not sit exactly at \(\varphi\). It sits at
\begin{equation}
\Delta_P:=P-\varphi.
\label{eq:deltaP}
\end{equation}

Informally: \(\Delta_P\) is the small amount by which the realized cell is displaced from exact self-similar equilibrium. That small displacement is what becomes the electromagnetic observation strength.

Why this step is needed: a fixed point needs a reference position. The detuning principle declares the golden ratio as the self-similar reference, while \(P-\varphi\) measures how far the actual cell sits away from it.

\section{Step 4: Boundary Gaussian Normalization}

The boundary normalization converts \(\Delta_P\) into the outside coupling readout:
\begin{equation}
\alpha_{\mathrm{ext}}(P):=\frac{P-\varphi}{\sqrt\pi}.
\label{eq:alphaext}
\end{equation}

Informally: the boundary Gaussian width normalizes the screen displacement. Dividing by \(\sqrt\pi\) turns the geometric displacement into a dimensionless observation strength. The Gaussian choice belongs to the declared detuning principle.

Why this step is needed: the outside screen variable is an area displacement. The inside electromagnetic variable is a coupling. The declared boundary normalization supplies the proposed conversion between those two readings.

\section{Step 5: The Inside Electromagnetic Readout}

Let
\begin{equation}
\ATh(P):=\alpha_{\mathrm{em}}^{-1}(0;P)
\label{eq:ATdef}
\end{equation}
be the inverse electromagnetic coupling at zero momentum returned by the
declared map at trial pixel \(P\). The inside coupling is
\begin{equation}
\alpha_{\mathrm{in}}(P):=\frac{1}{\ATh(P)}.
\label{eq:alphain}
\end{equation}

Informally: an observer inside the encoded world does not see \(P\) as a screen area. The observer sees the strength of electromagnetism. Since particle physicists quote the low-energy coupling as an inverse number near \(137\), the actual coupling is \(1/\ATh(P)\).

Why this step is needed: the measured fine-structure constant is the Thomson-limit coupling. The calculation must therefore end at zero momentum, after the electroweak anchor has been transported through all charged degrees of freedom.

\section{Step 6: The Fixed-Point Equation}

Conditional on the same-invariant bridge, closure requires the outside and inside couplings to agree:
\begin{equation}
\alpha_{\mathrm{ext}}(P)=\alpha_{\mathrm{in}}(P).
\label{eq:couplingmatch}
\end{equation}
Using Eqs.~\eqref{eq:alphaext} and \eqref{eq:alphain},
\begin{equation}
\frac{P-\varphi}{\sqrt\pi}
=
\frac{1}{\ATh(P)}.
\end{equation}
Equivalently,
\begin{equation}
\boxed{
H(P):=P-\varphi-\frac{\sqrt\pi}{\ATh(P)}=0
}
\label{eq:H}
\end{equation}
or
\begin{equation}
\boxed{
P=G(P):=\varphi+\frac{\sqrt\pi}{\ATh(P)}.
}
\label{eq:G}
\end{equation}

Informally: feed a trial pixel into the declared source chain. It returns a
Thomson endpoint. If the bridge proves that endpoint and the detuning are two
readings of one cell, closure selects a value that comes back unchanged.

Why this step is needed: a trial value is not enough. Once the bridge proves
that both readings belong to one cell, self-consistency requires them to
agree. That is why the conditional answer is a fixed point rather than a
one-way evaluation.

\section{Step 7: The Source Map}

For a trial \(P\), the source map first emits
\begin{align}
M_U(P)&=E_P\,\e^{-2\pi}\,P^{1/6},
\label{eq:MU}\\
E_{\mathrm{cell}}(P)&=\frac{E_P}{\sqrt P}.
\label{eq:Ecell}
\end{align}

Informally: \(M_U(P)\) is the high-scale source readout, and \(E_{\mathrm{cell}}(P)\) is the declared local-energy coordinate. Both are functions of the same trial input \(P\). Interpreting that input as one physical screen cell requires the same-cell bridge.

Why this step is needed: the inside electromagnetic coupling has to be generated from the same \(P\) used in the outside equation. These two formulas begin that inside calculation without adding a separate high-scale fit parameter.

\section{Step 8: Electroweak Transmutation}

The source branch uses
\begin{equation}
\beta_{\mathrm{EW}}=\Nc+1=4
\end{equation}
with \(\Nc=3\). The same value has a representation-theoretic reading in the
Standard Model gauge paper: the selected exterior matter package contains exactly four weak
doublets, \(3+1\) from the quark and lepton sectors of one generation, so
\(\Nc+1\) and the weak-doublet multiplicity name one coefficient. The
identification of that abstract load with the physical screen load is the
separate common-carrier receipt. The branch defines
\begin{equation}
v(P,\alpha_U)
=
E_{\mathrm{cell}}(P)
\exp\!\left[-\frac{2\pi}{\beta_{\mathrm{EW}}\alpha_U}\right].
\label{eq:v}
\end{equation}

Informally: the weak scale is exponentially lower than the cell scale. The unified coupling controls that descent, so changing \(\alpha_U\) changes the electroweak scale sharply.

Why this step is needed: the fine-structure constant is measured at low energy, and the source map begins at the cell scale. The transmutation formula explains how the electroweak scale is produced from the same source data. After independent direct public-record closure, a positive refinement-natural identification of the screen and electroweak load carriers gives the conditional bridge \(N_{\rm bridge}=\pi\exp[6\pi/(P\alpha_U)]\), certified by interval arithmetic at \(3.53\times10^{122}\). The frozen FZ-01 capacity registration carries this uncorrected bridge value, \(3.532131543\times10^{122}\), as the null alternative to its band \([3.30836,3.31137]\times10^{122}\) \cite{ophladder}. The bridge does not construct the public-record readback.

\section{Step 9: One-Loop Gauge Running}

For \(i=1,2,3\), define
\begin{equation}
(b_1,b_2,b_3)=\left(\frac{33}{5},\,1,\,-3\right),
\label{eq:b}
\end{equation}
and run the couplings by
\begin{equation}
\alpha_i^{-1}(\mu;P,\alpha_U)
=
\alpha_U^{-1}
+
\frac{b_i}{2\pi}\log\!\left(\frac{M_U(P)}{\mu}\right).
\label{eq:running}
\end{equation}
The triple \((33/5,1,-3)\) is the one-loop MSSM beta-coefficient set. Its use
is a declared structural selection: the Standard Model coefficient triple has
no pixel-residual root in the declared scan, no supersymmetric threshold
matching enters the map, and the convention is a counted structural selection.

Informally: once \(\alpha_U\) is guessed, all three gauge couplings at lower scales are fixed. The same high-scale coupling determines \(\alpha_1\), \(\alpha_2\), and \(\alpha_3\).

Why this step is needed: the electromagnetic coupling is a mixture of the weak and hypercharge couplings. Running the gauge couplings supplies the ingredients needed to form that mixture at the \(Z\)-scale.

\section{Step 10: The Self-Consistent \(Z\)-Scale}

The hypercharge coupling in Standard Model normalization is
\begin{equation}
\alpha_Y(\mu;P)=\frac{3}{5}\alpha_1(\mu;P).
\label{eq:alphaY}
\end{equation}
The tree-level \(Z\)-mass readout is
\begin{equation}
m_Z(\mu;P,\alpha_U)
=
\frac{v(P,\alpha_U)}{2}
\sqrt{4\pi\alpha_2(\mu;P,\alpha_U)+4\pi\alpha_Y(\mu;P,\alpha_U)}.
\label{eq:mZtree}
\end{equation}
The source point uses the self-consistency condition
\begin{equation}
\mu=m_Z(\mu;P,\alpha_U).
\label{eq:mZfixed}
\end{equation}

Informally: the \(Z\)-scale is determined inside the calculation. It is the scale at which the running couplings and the electroweak transmutation formula reproduce their own \(Z\)-mass readout.

Why this step is needed: using a measured \(m_Z\) here would hide experimental input inside the source map. The self-consistency condition keeps the source calculation closed.

\section{Step 11: Heat-Kernel Gauge Closure}

For \(\SU(2)\), the irreducible representations are labeled by \(j=n/2\), \(n=0,1,\ldots,N_2\), with
\begin{equation}
d_j=2j+1,
\qquad
C_2(j)=j(j+1).
\label{eq:su2rep}
\end{equation}
For \(\SU(3)\), representations are labeled by highest weights \((p,q)\), \(0\le p,q\le N_3\), with
\begin{align}
d_{p,q}&=\frac{(p+1)(q+1)(p+q+2)}{2},
\label{eq:su3dim}\\
C_2(p,q)&=\frac{p^2+q^2+pq+3p+3q}{3}.
\label{eq:su3casimir}
\end{align}

For a compact group \(G\) in this finite representation cutoff, define
\begin{align}
Z_G(t)&=\sum_R d_R\,\e^{-tC_2(R)},\\
\bar\ell_G(t)&=
\frac{1}{Z_G(t)}
\sum_R d_R\,\e^{-tC_2(R)}\log d_R.
\label{eq:ellbar}
\end{align}
The displayed evaluations use the representation cutoffs \(N_2=120\) and
\(N_3=90\). At the closure arguments below, \(t_2\approx1.3\) and
\(t_3\approx4.7\), the first omitted terms carry weights below
\(\e^{-4000}\), so the cutoffs are inert at every displayed precision, and
the interval certificate bounds the tails beyond any finite cutoff.

The heat-kernel source parameters are
\begin{equation}
t_2=4\pi^2\alpha_2(m_Z;P,\alpha_U),
\qquad
t_3=4\pi^2\alpha_3(m_Z;P,\alpha_U).
\label{eq:t23}
\end{equation}
The pixel-closure equation is
\begin{equation}
\boxed{
\bar\ell_{\SU(2)}(t_2)+\bar\ell_{\SU(3)}(t_3)=\frac{P}{4}.
}
\label{eq:pixelclosure}
\end{equation}
For a trial \(P\), Eq.~\eqref{eq:pixelclosure} is solved for \(\alpha_U(P)\).
The closure sums the \(\SU(2)\) and \(\SU(3)\) representation towers only; the
selection that hypercharge carries no representation entropy against the
\(P/4\) budget is a declared choice. The repair round count used by the
finite current construction counts the \(\mathfrak u(1)\) generator
(\(8+3+1=12\)); the two
conventions are distinct declared choices.

Informally: the representation entropy carried by the weak and color sectors must match the pixel capacity assigned to the cell. This step is what locks the unified coupling to the same \(P\) that appears in the outer equation.

Why this step is needed: \(\alpha_U\) should not be chosen by hand. The heat-kernel closure turns the finite representation content of the gauge sectors into an equation for \(\alpha_U(P)\).

\section{Step 12: The Electroweak Anchor}

Once \(\alpha_U(P)\) and \(m_Z(P)\) are solved, define
\begin{align}
\alpha_Y(m_Z;P)&=\frac{3}{5}\alpha_1(m_Z;P),\\
\alpha_{\mathrm{em}}(m_Z^2;P)
&=
\left(
\frac{1}{\alpha_2(m_Z;P)}
+\frac{1}{\alpha_Y(m_Z;P)}
\right)^{-1}.
\label{eq:alphaem}
\end{align}
The source-locked electroweak anchor is
\begin{equation}
\boxed{
\AZ(P):=\alpha_{\mathrm{em}}^{-1}(m_Z^2;P).
}
\label{eq:AZ}
\end{equation}
The weak mixing readout is
\begin{equation}
\sin^2\theta_W(m_Z;P)=
\frac{\alpha_{\mathrm{em}}(m_Z^2;P)}{\alpha_2(m_Z;P)}.
\label{eq:sin2}
\end{equation}

Informally: the source map has reached the electromagnetic coupling at the electroweak anchor scale. The low-energy \(1/137\) number appears only after this anchor is transported to zero momentum.

Why this step is needed: \(\AZ(P)\) is the clean place where the source map meets standard electroweak physics. From here the problem becomes a transport problem for the electromagnetic current.

\section{Step 13: Frozen Charged-Spectrum Continuation Used by the Transport}

The one-loop kernel below is exact for supplied fermion masses, but the charged
spectrum supplied to it in this section is a phenomenological continuation,
not an OPH-derived mass theorem.  It uses \(\Nc=3\), \(\Ng=3\), and the
additional assumptions
\begin{equation}
\epsilon=\frac16,
\qquad
\delta=\frac{\beta_{\mathrm{EW}}}{2\Nc\Ng}
=\frac{2}{9}.
\label{eq:epsdelta}
\end{equation}
No physical charged-spectrum derivation is claimed here; such a derivation
requires a source-to-mass attachment, a phase-transport law fixing \(\delta=2/9\),
and an attachment of the ordered roots to physical charged-family lines. For a
positive-semidefinite square-root-mass carrier
\[
C=aI+\rho(e^{i\delta}R+e^{-i\delta}R^2),
\qquad
Q=\frac{1+2(\rho/a)^2}{3},
\]
so \(\rho/a=1/\sqrt2\) is exactly equivalent to the empirical Koide value
\(Q=2/3\). At balance this identity holds throughout the positive
chamber \(|\delta|\leq\pi/12\pmod{2\pi/3}\), and it does not select the phase.
The balance is a target-informed premise and the phase \(\delta\) is
supplied. The finite tracial construction that realizes the balance, and
the maps a physical charged-lepton reading would additionally require, are
stated in the Koide paper~\cite{oph-koide}.
Define three Koide roots
\begin{equation}
r_k=1+\sqrt2\cos\!\left(\delta+\frac{2\pi k}{3}\right),
\qquad k=0,1,2,
\label{eq:koide}
\end{equation}
then sort them in increasing order and write the sorted list as
\((r_1,r_2,r_3)\).

The quark exponent vectors are
\begin{equation}
\bm n_u=(2\Nc,\Nc,0)=(6,3,0),
\qquad
\bm n_d=(2\Nc,\Nc+1,\Nc-1)=(6,4,2).
\label{eq:quarkexponents}
\end{equation}
With \(v=v(P)\),
\begin{align}
m_u&=\frac{v}{\sqrt2}\epsilon^6,
&
m_c&=\frac{v}{\sqrt2}\epsilon^3,
&
m_t&=\frac{v}{\sqrt2},\\
m_d&=\frac{v}{\sqrt2}\epsilon^6,
&
m_s&=\frac{v}{\sqrt2}\epsilon^4,
&
m_b&=\frac{v}{\sqrt2}\epsilon^2.
\label{eq:quarkmasses}
\end{align}

For charged leptons the exponent vector is
\begin{equation}
\bm n_e=(7,4,3).
\label{eq:lepexp}
\end{equation}
This vector and the normalization below are declared phenomenological transport
inputs. The declared premises do not derive
\(\bm n_e=(7,4,3)\) or the resulting determinant identity
\(\det M_e=v(P)^3/(2\,6^{14})\). The electron entry is selected by a
finite scan: with \(n_\tau=N_g=3\) and \(n_\mu=N_g+1=4\) fixed by the
generation count, the electron exponent runs over the integers \(5\) to
\(11\), and the entry with the smallest maximal pairwise logarithmic residual
between the squared balanced-circulant roots at
\(\delta=(N_c+1)/(2N_cN_g)\) and the corresponding powers of \(\epsilon=1/6\)
is kept; the minimizer is \(7\). The scan compares the declared circulant
roots with the \(\epsilon\) ladder and reads no measured mass, and the
vector is a declared transport input all the same. Every common shift
\((7+k,4+k,3+k)\) has the same
ratio residual while changing the determinant.  The extra factor \(2^{1/6}\)
below is asserted rather than derived.
Define
\begin{equation}
\log g_c
=
\frac{1}{3}
\sum_{a=1}^3
\log\!\left(
\frac{r_a^2\sqrt2\,6^{n_{e,a}}}{v}
\right),
\qquad
s_0=\e^{-\log g_c},
\qquad
s_e=s_0\,2^{1/6}.
\label{eq:lepscale}
\end{equation}
Then
\begin{equation}
m_e=s_e r_1^2,
\qquad
m_\mu=s_e r_2^2,
\qquad
m_\tau=s_e r_3^2.
\label{eq:lepmasses}
\end{equation}

Equations~\eqref{eq:lepscale}--\eqref{eq:lepmasses} define a transport input,
not a source-derived charged-lepton mass law. No charged RG/threshold theorem
maps this coordinate to physical pole masses.

Informally: this section specifies a phenomenological spectrum used by the
transport diagnostic.  The lepton vacuum-polarization kernel is an exact
one-loop calculation conditional on those masses; that exactness does not
determine the masses themselves.  The quark part is also a perturbative
continuation. The confined hadronic QCD spectral measure is a separate
low-energy object.

Why this step is needed: vacuum polarization depends on charged particles. The
transport diagnostic cannot be evaluated until a charged spectrum and its
charges are supplied.  A pure OPH source theorem would first have to replace
the continuation spectrum by independently derived charged masses.

\section{Step 14: Exact One-Loop Fermion Transport}

For a fermion of mass \(m_f\), electric charge \(Q_f\), and multiplicity \(N_f\), define
\begin{equation}
K_f(Q^2;m_f,Q_f,N_f)
=
\frac{2N_fQ_f^2}{\pi}
\int_0^1 x(1-x)
\log\!\left(1+\frac{Q^2x(1-x)}{m_f^2}\right)\dd x.
\label{eq:kernelint}
\end{equation}

The integral has the following closed form. Let
\begin{equation}
z=\frac{Q^2}{m_f^2},
\qquad
a=\frac{z}{4}.
\end{equation}
Then
\begin{equation}
\int_0^1 x(1-x)\log(1+zx(1-x))\dd x
=
-\frac{5}{18}
+\frac{1}{6a}
+\frac{(2a-1)\sqrt{1+a}\,\operatorname{asinh}(\sqrt a)}
{6a^{3/2}},
\label{eq:kernelclosed}
\end{equation}
with
\begin{equation}
\operatorname{asinh}(\sqrt a)
=
\log(\sqrt a+\sqrt{1+a}).
\end{equation}

The lepton contribution is
\begin{equation}
\Dlep(P)=
K_e(m_Z(P)^2;m_e,1,1)
+K_\mu(m_Z(P)^2;m_\mu,1,1)
+K_\tau(m_Z(P)^2;m_\tau,1,1).
\label{eq:Dlep}
\end{equation}
The naive five-quark contribution is
\begin{align}
\Delta_q^{\mathrm{naive}}(P)
&=
K_u(m_Z(P)^2;m_u,2/3,3)
+K_d(m_Z(P)^2;m_d,-1/3,3)
\nonumber\\
&\quad
+K_s(m_Z(P)^2;m_s,-1/3,3)
+K_c(m_Z(P)^2;m_c,2/3,3)
+K_b(m_Z(P)^2;m_b,-1/3,3).
\label{eq:Dqnaive}
\end{align}
The declared hadronic screening ansatz is
\begin{equation}
S_{\mathrm{calc}}(P)=
1-\frac{\Nc\alpha_3(m_Z;P)}{\pi}.
\label{eq:Scalc}
\end{equation}
The form of \(S_{\mathrm{calc}}\) is a declared selection of the quark
continuation and carries no perturbative derivation: the leading perturbative
correction to quark vacuum polarization is the enhancement
\(R\to R_0(1+\alpha_s/\pi+\cdots)\) with no \(\Nc\) factor, about \(+3.7\)
percent at \(m_Z\), while \(S_{\mathrm{calc}}\approx0.887\) is an \(11\)
percent suppression standing in for confinement cutting off the light-quark
logarithms.
Thus
\begin{equation}
\Dcalc(P)
=
\Dlep(P)+S_{\mathrm{calc}}(P)\Delta_q^{\mathrm{naive}}(P).
\label{eq:Dcalc}
\end{equation}

Informally: the supplied charged-lepton continuation is propagated through an
exact one-loop kernel. Quarks are treated by the one-loop kernel times the
declared screening ansatz. The expected gap is the confined hadronic spectral
transport plus the matching terms needed to use the same electromagnetic-current convention all the way to the endpoint.

Why this step is needed: the Thomson endpoint is not the same as the \(Z\)-scale anchor. Charged particles screen the electromagnetic current between those scales, and the kernel is the mathematical form of that screening.

\section{Step 15: The Source Calculation Endpoint}

The calculated endpoint is
\begin{equation}
A_{\mathrm{calc}}(P)
=
\AZ(P)+\Dcalc(P).
\label{eq:Acalc}
\end{equation}
Putting \(A_{\mathrm{calc}}\) into Eq.~\eqref{eq:G} gives the source map
\begin{equation}
G_{\mathrm{calc}}(P)
=
\varphi+\frac{\sqrt\pi}{A_{\mathrm{calc}}(P)}.
\label{eq:Gcalc}
\end{equation}
The numerical solve uses
\begin{equation}
G_{\mathrm{calc}}(P_{\mathrm{root}})=P_{\mathrm{root}}
\label{eq:Proot}
\end{equation}
and emits
\begin{align}
P_{\mathrm{root}}
&=
1.630972095858897376964513903506955628479\ldots,
\label{eq:Prootnum}\\
\alphainv_{\mathrm{root}}
&=
136.9948351774129372952894294644369028576\ldots.
\label{eq:alpharoot}
\end{align}
The fixed point is interval-certified as
existing and unique on its interval (interval certificate,
available in the public artifact bundle \cite{oph-code};
enclosure widths \(6.8\times10^{-28}\) in \(P\) and \(7.2\times10^{-24}\) in
\(\alpha^{-1}\), Lipschitz bound \(L\le0.0724\)).
The source anchor at the certified root point is
\begin{equation}
\AZ(P_{\mathrm{root}})
=
128.3082680579875973479040576140847588750331\ldots.
\label{eq:AZroot}
\end{equation}
The calculated transport contribution is
\begin{equation}
\Dcalc(P_{\mathrm{root}})
=
8.68656711942533994738537185035214398252\ldots.
\label{eq:Dcalcroot}
\end{equation}

Informally: the fixed-point algebra lands at a stable value near \(137\), with no measured alpha inserted into the solve. Its closeness to the measured value depends on the declared screening ansatz; with the measured hadronic payload in its place the endpoint is \(136.3827548\) (Section~\ref{sec:pdg-step}). The numerical interval does not bound the missing physical contributions.

Why this step is needed: this is the non-circular check. It shows what the OPH source chain produces before the calibrated hadronic endpoint correction is added.

The P-dependent detector of the executable observer construction takes this root as its closure
input. Under a chosen charge assignment it sets
\(g(P_{\mathrm{root}})=\sqrt{4\pi\alpha(P_{\mathrm{root}})}=0.3028676131\ldots\) and records a
binary detector observation whose probabilities depend on \(P\) and which feeds the next native
repair; that use leaves the endpoint status of the root unchanged \cite{ophpresponse}.

\section{Step 16: The CODATA/NIST Measured Endpoint and Comparison Pixel}

The CODATA/NIST 2022 inverse fine-structure constant \cite{nist-value,nist-wall} is
\begin{equation}
A_{\mathrm C}=137.035999177,
\qquad
\sigma_A=0.000000021,
\label{eq:codata}
\end{equation}
with concise form \(137.035999177(21)\). The corresponding comparison pixel is
\begin{equation}
\Ppub
=
\varphi+\frac{\sqrt\pi}{A_{\mathrm C}}
=
\decnum{1.6309682094039593248792798477826489413359828516279250606661507533907793398933432}
\label{eq:Ppub}
\end{equation}
The CODATA/NIST central coupling is
\begin{equation}
\alpha(0)
=
\frac{1}{A_{\mathrm C}}
=
\decnum{0.0072973525643314250302457952646916832280660213133653604957798803819933561573639928}
\label{eq:alphapub}
\end{equation}

Informally: once the measured Thomson endpoint is supplied, the outer OPH equation fixes the corresponding comparison pixel immediately. It is the direct readout of one equation.

Why this step is needed: the measured endpoint and the comparison pixel are two algebraic forms of the same comparison definition. Reporting both makes the measured provenance of the pixel explicit; it does not identify that coordinate with a source-selected screen cell.

\paragraph{Two-\(P\) provenance audit.}
There are two distinct pixel numerals in this paper.  The source-only
solve gives \(P_{\mathrm{root}}\) in Eq.~\eqref{eq:Prootnum}; it is the
branch that audits what the OPH source chain emits before a same-scheme
hadronic endpoint correction is supplied.  The comparison value \(\Ppub\) in
Eq.~\eqref{eq:Ppub} is different by about \(3.9\times10^{-6}\) in \(P\)-space
because it is defined from the measured CODATA/NIST Thomson endpoint.  Thus
the published comparison pixel is a measured-endpoint display, not a
derivation of \(\alpha\).  Downstream quantities that use \(\Ppub\) inherit
that measured endpoint by definition; downstream quantities that use
\(P_{\mathrm{root}}\) are source-branch diagnostics. The registered conditional primordial tilt
\(n_s=1-P/48=0.9660214956\ldots\) is one such downstream quantity; it is evaluated at \(\Ppub\)
\cite{ophpublicdata,ophladder}.

\section{Step 17: Endpoint Accounting at the CODATA/NIST Comparison Pixel}

At \(P=\Ppub\), the source point gives
\begin{align}
\AZ(\Ppub)
&=
\decnum{128.30796547328624820996110874175671618724547618036535646005342169635117784168285644078724728}
\label{eq:AZpub}\\
\Dcalc(\Ppub)
&=
\decnum{8.6865678427085284009854425428859697682672217376487364233784577389993784459106783080961179590}
\label{eq:Dcalcpub}
\end{align}
These are non-certified artifact displays: the quoted digit strings are the
run's output at its declared convergence budget, and independent
recomputation agrees with \(\AZ(\Ppub)\) to 27 significant digits rather
than to the full displayed precision; the certified interval artifacts are
unaffected and land inside their intervals far deeper.
The transport required by the CODATA/NIST measured endpoint is
\begin{equation}
\Dreq(\Ppub)
=
A_{\mathrm C}-\AZ(\Ppub)
=
\decnum{8.72803370371375179003889125824328381275452381963464353994657830364882215831714355921275272}
\label{eq:Dreqpub}
\end{equation}
Therefore the source-side residual at the CODATA/NIST comparison pixel is
\begin{align}
\Rq(\Ppub)
&=
\Dreq(\Ppub)-\Dcalc(\Ppub)
\nonumber\\
&=
\decnum{0.04146586100522338905344871535731404448730208198590711656812056464944371240646525111663476}
\label{eq:RQpub}
\end{align}

Informally: the CODATA/NIST measured value differs from the pure source calculation by a small, precisely located inverse-alpha amount. Neither declared map computes same-scheme low-energy hadronic transport to the Thomson endpoint, and inserting the measured hadronic payload widens the difference rather than closing it (Section~\ref{sec:pdg-step}). The fixed-point solve and this accounting do not identify the physical source of the gap \cite{ophpredictions,ophpostdiction}.

Why this step is needed: this accounting prevents the hadronic contribution from being hidden inside the final number. It shows the source anchor, the calculated transport, the required transport, and the residual in the same units.

\section{Step 18: The OPH Output and the Required Hadronic Correction}

The fixed-point calculation emits
\begin{equation}
A_{\mathrm{calc}}(P_{\mathrm{root}})
=
\alphainv_{\mathrm{root}}
=
\decnum{136.994835177412937295289429464436902857561206151035393184502}.
\label{eq:Acalcfp}
\end{equation}
The interval certificate and continuous-integration identity test are included
in the public artifact bundle \cite{oph-code}.
At the CODATA-derived comparison pixel, the finite-screen unified gauge-width contribution is
\[
\alpha_U(\Ppub)=0.041124336195630495.
\]
Combining that comparison-pixel gauge width with Equation~\eqref{eq:Acalcfp}
gives the mixed-provenance
no-hadron diagnostic
\begin{equation}
A_{\alpha_U}^{\mathrm{fp}}
=A_{\mathrm{calc}}(P_{\mathrm{root}})+\alpha_U(\Ppub)
=137.035959513608567790289429464437\ldots,
\label{eq:barealphaUaddition}
\end{equation}
or
\begin{equation}
\alpha_{\mathrm{src}+U}^{\mathrm{fp}}
=\bigl(A_{\alpha_U}^{\mathrm{fp}}\bigr)^{-1}
=0.00729735467646135213287275959657\ldots.
\label{eq:barealphaUalpha}
\end{equation}
This \(A_{\alpha_U}^{\mathrm{fp}}\) value is comparison bookkeeping, not a source-only prediction:
\(\alpha_{\mathrm{root}}^{-1}\) is evaluated on the source root, while \(\alpha_U(\Ppub)\) inherits
the measured endpoint through \(\Ppub\), and the sum is not a fixed point of any single declared map.

Compared with the CODATA/NIST central measured endpoint,
\begin{equation}
A_{\mathrm C}=137.035999177,
\qquad
\sigma_A=0.000000021,
\end{equation}
the CODATA-minus-diagnostic bookkeeping gap is
\begin{equation}
\Dhadreq^{\mathrm{fp}}
=
A_{\mathrm C}-A_{\alpha_U}^{\mathrm{fp}}
=
\decnum{0.000039663391432209710570535563}\ldots.
\label{eq:Dempfp}
\end{equation}
Equivalently,
\begin{equation}
\Dhadreq^{\mathrm{fp}}
=0.0000396634\pm0.000000021,
\qquad
\frac{\Dhadreq^{\mathrm{fp}}}{A_{\mathrm C}}
=2.89438\times10^{-7}.
\label{eq:Dhadrequncertainty}
\end{equation}
This is the inverse-alpha gap between the mixed diagnostic and the CODATA/NIST central value. It is
the gap left by the declared screening ansatz; inserting the measured hadronic vacuum
polarization instead does not close it (Section~\ref{sec:pdg-step}). The subtraction neither
calculates same-scheme transport of the Ward-projected electromagnetic current nor
identifies the physical source of the gap.
The \(2.89438\times10^{-7}\) ratio describes the mixed diagnostic only. The distinct
self-consistent gauge-width map has fixed point
\(\alpha^{-1}=137.035660136946577\ldots\) and relative residual
\(2.5\times10^{-6}\), about \(1.6\times10^{4}\) measurement sigma.

For bookkeeping, the full root-to-CODATA difference is
\begin{align}
\Dcalib^{\mathrm{fp}}
&=
A_{\mathrm C}-A_{\mathrm{calc}}(P_{\mathrm{root}})
\nonumber\\
&=
\decnum{0.041163999587062704710570535563097142439}
\nonumber\\
&=
\alpha_U(\Ppub)+\Dhadreq^{\mathrm{fp}}.
\label{eq:rootcodatafp}
\end{align}
This larger difference is bookkeeping: it contains the comparison-pixel unified gauge-width
contribution and the CODATA-minus-diagnostic gap.

Equivalently, the CODATA/NIST central value can be written as
\begin{equation}
A_{\mathrm C}
=
\alpha_{\mathrm{root}}^{-1}+\alpha_U(\Ppub)\,C_{24,Q},
\qquad
C_{24,Q}=1.0009644749338573880401357114045152\ldots.
\label{eq:c24qfactor}
\end{equation}
The factor \(C_{24,Q}\) is a back-solved calibration: compact comparison
accounting, solved from the measured endpoint, for the finite same-scheme hadronic endpoint
transport needed for source-only closure. Equivalently,
\(\Dhadreq^{\mathrm{fp}}=\alpha_U(\Ppub)(C_{24,Q}-1)\). The factor is not a source-emitted
hadronic/QCD payload.

Here \(A_{\mathrm C}\) is the CODATA/NIST 2022 inverse fine-structure constant \cite{nist-value,nist-2022}. Closing the endpoint gap requires a source-derived hadronic spectral calculation, which this paper does not supply.

The endpoint decomposition is
\begin{equation}
\ATh(P)
=
\AZ(P)+\Dlep(P)+\Dhad(P)+\DEW(P).
\label{eq:ATsplit}
\end{equation}
The calculated transport contains the exact one-loop charged-lepton contribution and a structured quark-screening continuation. The residual pure source object can be written compactly as
\begin{equation}
\Rq(P)
=
\Dhad(P)+\DEW(P)
-\bigl[\Dcalc(P)-\Dlep(P)\bigr],
\label{eq:RQmeaning}
\end{equation}
where the right side is understood in the same renormalization and endpoint scheme as \(\AZ(P)\).

In the screening notation used by the endpoint calculation, let
\begin{equation}
x(P)=\frac{\Nc\alpha_3(m_Z;P)}{\pi}.
\label{eq:xqcd}
\end{equation}
At \(P=\Ppub\),
\begin{align}
S_{\mathrm{required}}
&=
\decnum{0.89540013264765879780580028318167064130770986481229164844830591011545242521273086888107864576}
\label{eq:Sreq}\\
c_Q
&=
\decnum{0.65802575992715543563823017023236005042492009907058639608566007832470711257322011342309305088}
\label{eq:cQ}
\end{align}
where
\begin{equation}
S_{\mathrm{required}}=1-x+c_Qx^2.
\label{eq:cQdef}
\end{equation}

Informally: the needed hadronic and endpoint correction can be summarized as a small second-order screening coefficient. A pure source proof has to emit this coefficient from a Ward-projected hadronic spectral measure.

Why this step is needed: hadrons are not pointlike free quarks at low energy. The coefficient records the screening that the CODATA/NIST endpoint requires in this endpoint convention; it is back-solved from that endpoint.

\section{Step 19: Why Raw PDG \(\Delta\alpha_{\mathrm{had}}^{(5)}(M_Z)\) Is a Different Quantity}
\label{sec:pdg-step}
\label{sec:raw-pdg}

The direct PDG/CERN diagnostic uses
\begin{align}
\AZ(\Ppub)&=
\decnum{128.30796547328624820996110874175671618724547618036535646005342169635117784168285644}\\
\Dlep(\Ppub)&=
\decnum{4.3093978664522040271317438975344894018487156605576773194711528089665680313257906466129}\\
\Delta\alpha_{\mathrm{had}}^{(5)}(M_Z)&=0.02761.
\end{align}
It forms
\begin{equation}
A_L=\AZ+\Dlep,
\qquad
A_{\mathrm{PDGdiag}}=\frac{A_L}{1-\Delta\alpha_{\mathrm{had}}^{(5)}(M_Z)}.
\label{eq:pdgdiag}
\end{equation}
Numerically,
\begin{equation}
A_{\mathrm{PDGdiag}}
=
\decnum{136.382895072695577121415124218977165118002233508081154453999500720202537945689123794581289}
\label{eq:pdgdiagval}
\end{equation}
This differs from \(137.035999177\). Holding this raw hadronic denominator shift fixed, the target would require
\begin{equation}
\Delta\alpha_{\mathrm{had,req}}
=
\decnum{0.032244343557887288822268499236957947422007347612226081506593089221871380926117950141438549}
\label{eq:pdgreq}
\end{equation}
or a same-scheme source-anchor bridge of
\begin{equation}
\begin{gathered}
\bigl(\Delta\alpha_{\mathrm{had,req}}-\Delta\alpha_{\mathrm{had}}^{(5)}(M_Z)\bigr)\cdot\alphainv(0)
=\\
\decnum{0.6350718999845777629071473607087944109058081590769662204754254946822541269913529133871}
\end{gathered}
\label{eq:bridgegap}
\end{equation}
inverse-alpha units, with \(\Delta\alpha_{\mathrm{had,req}}\) from
Eq.~\eqref{eq:pdgreq}, \(\Delta\alpha_{\mathrm{had}}^{(5)}(M_Z)=0.02761\), and
\(\alphainv(0)=137.035999177\) the CODATA central endpoint.
With the pinned KNT19 payload \(\Delta\alpha_{\mathrm{had}}^{(5)}(M_Z^2)=0.027609\pm0.000112\),
the same construction gives \(\alpha^{-1}(0)=136.3827548\), inside the interval from \(136.3670481\) to \(136.3984652\).
The measured value lies outside this interval. The corresponding same-scheme anchor gap lies in
\([0.6199,0.6506]\) inverse-alpha units and contains the standard on-shell reference deficit
\(0.6309\); both numbers measure the same anchor gap, computed once with the OPH lepton transport
and once with standard running, so their agreement checks the bookkeeping and does not make the
endpoint compatible with the measured value. This compare-only row does not identify the
physical source of the gap \cite{ophpostdiction}.

Informally: the raw electroweak-review hadronic running number is a useful diagnostic. The OPH endpoint needs the same electromagnetic current and the same inverse-alpha endpoint convention used by \(\AZ(P)\). Mixing the two conventions moves the answer by a visible amount.

Why this step is needed: a common mistake is to insert a published \(\Delta\alpha_{\mathrm{had}}^{(5)}(M_Z)\) number as though it were the OPH endpoint contribution. This section shows the numerical consequence of that convention mismatch.

\section{Step 20: The Measured Hadronic Spectral Input}

The standard measured hadronic vacuum-polarization relation uses
\begin{equation}
\Delta\alpha_{\mathrm{had}}(q^2)
=
-\frac{\alpha q^2}{3\pi}\,
\mathrm{P.V.}\!\int \frac{R(s)}{s(s-q^2)}\dd s,
\label{eq:dispersion}
\end{equation}
where \(R(s)\) is the measured ratio of the bare \(e^+e^-\to\mathrm{hadrons}\) cross section to the pointlike \(e^+e^-\to\mu^+\mu^-\) cross section.

The independently documented measured quantity in the standard electroweak literature is usually quoted as \(\Delta\alpha_{\mathrm{had}}^{(5)}(M_Z^2)\). It is a dimensionless change in the running electromagnetic coupling at the \(Z\)-boson mass, not the inverse-alpha correction \(\Dhadreq^{\mathrm{fp}}\) used in Eq.~\eqref{eq:Dempfp}. For example, Davier, Hoecker, Malaescu, and Zhang report
\[
\Delta\alpha_{\mathrm{had}}^{(5)}(M_Z^2)=(275.7\pm1.0)\times10^{-4}
\]
from \(e^+e^-\)-based data \cite{dhmz}. A perturbative update by Erler and Ferro-Hernandez gives
\[
\Delta\alpha_{\mathrm{had}}^{(5)}(M_Z^2)
=\left(276.29\pm0.38\pm0.62\right)\times10^{-4}
\]
when low-energy cross-section data are used as input \cite{erler-ferro}.

The comparison residual in Eq.~\eqref{eq:Dempfp} is different. It is expressed in inverse-alpha units and in the endpoint convention used by \(\AZ(P)\). The standard published hadronic-running values document the physics source of the correction. They do not by themselves give the OPH same-scheme inverse-alpha correction to arbitrary precision. The community does publish the data machinery needed for this kind of calculation. Jegerlehner's alphaQED package provides hadronic running routines, covariance data, and an integration routine for custom kernels \cite{alphaqed}. That is the right empirical route once the OPH endpoint kernel and finite same-scheme remainder are fixed.

In the OPH endpoint convention this is represented as a same-current spectral functional,
\begin{equation}
\Dhad(P)
=
\frac{m_Z(P)^2}{3\pi}
\int \frac{\rho_Q(s;P)}{s[s+m_Z(P)^2]}\dd s,
\label{eq:ophhadron}
\end{equation}
plus the same-scheme finite remainder needed by Eq.~\eqref{eq:ATsplit}.

Informally: the measured hadronic input supplies the electromagnetic spectral information required for the endpoint. The single published number \(0.0276\) is one weighted integral of that spectral information. The OPH endpoint needs a different weighted integral plus the same-scheme finite remainder. A pure source theorem requires the same spectral object from OPH source data.

Why this step is needed: the dispersion relation explains why measured \(e^+e^-\to\mathrm{hadrons}\) data are the right empirical input. They measure the electromagnetic spectral function that the endpoint transport requires.

\section{No-Go: The 24-Slot Register Scale Does Not Determine the Hadronic Residual}

\begin{theorem}[Coarse 24-slot register data do not determine the Thomson residual]
The specified OPH source objects \(P\), \(\alpha_U(P)\), \(N_c=3\), \(N_g=3\), and the
24-slot repair-register count do not determine the exact low-energy Thomson residual.
\end{theorem}

\begin{proof}
The hadronic endpoint contribution depends on the Ward-projected spectral functional
\begin{equation}
\Delta_{\mathrm{had}}(P)
=
\frac{m_Z(P)^2}{3\pi}
\int_0^\infty
\frac{\rho_Q(s;P)}{s[s+m_Z(P)^2]}\dd s
\label{eq:hadronkernelnogo}
\end{equation}
up to the same-scheme finite endpoint remainder. Let
\[
K_P(s)=\frac{m_Z(P)^2}{3\pi s[s+m_Z(P)^2]}.
\]
This kernel is not constant. Two positive spectral measures can therefore have the same coarse
register count, the same leading normalization, and the same electroweak branch data while giving
different residuals. For example, with \(s_1\ne s_2\),
\[
\rho_1(s)=M\delta(s-s_1),
\qquad
\rho_2(s)=M\delta(s-s_2)
\]
have equal total mass, but
\[
\int K_P(s)\rho_1(s)\dd s=M K_P(s_1),
\qquad
\int K_P(s)\rho_2(s)\dd s=M K_P(s_2),
\]
which are unequal because \(K_P\) is strictly decreasing on positive support.
Thus the register count does not select the spectral measure or the residual.
Even at fixed total mass \(M\), moving a positive atom toward zero or toward
infinity makes its moment arbitrarily large or small. A quantitative scale
bound needs additional support or moment assumptions.
\end{proof}

The hadronic precision audit sharpens the no-go statement. The
\(\rho_Q(s;P)\) object used here is the two-current marginal of a larger source-derived hadronic spectral
backend: a source QCD quotient ensemble, source parameter map, Ward-normalized current
data, Stieltjes/Jacobi export, same-scheme remainder, and systematics/no-target-leak receipt
bundle. That marginal can close the running-\(\alpha\)/HVP transport only after it is emitted by
the backend. It cannot be recycled as a full hadronic precision source, because HLbL and rare
decay long-distance amplitudes require four-current and transition spectral data from the same
source law.

\begin{theorem}[Tautological 24-Jacobi realization]
For any positive residual \(R>0\), there is an exact 24-point Jacobi/Stieltjes representation that
reproduces \(R\). Such a representation is non-predictive unless its nodes, weights,
normalization, and finite remainder are emitted by OPH source rules before comparison with the
Thomson endpoint.
\end{theorem}

\begin{proof}
Fix \(m>0\). Choose pairwise distinct positive spectral nodes
\(s_1,\ldots,s_{24}>0\) and positive normalized weights
\(u_j>0\) with \(\sum_j u_j=1\). The finite inverse Stieltjes/Favard construction gives a
positive \(24\times24\) Jacobi matrix \(J_{24,Q}\) whose \(e_1\)-spectral measure is
\[
\dd\mu_0(s)=\sum_{j=1}^{24}u_j\delta(s-s_j).
\]
Equivalently,
\[
e_1^\top f(J_{24,Q})e_1=\sum_{j=1}^{24}u_j f(s_j)
\]
for every function \(f\) on the spectrum. Define
\[
D(J_{24,Q};m)
=
e_1^\top\!\left[J_{24,Q}^{-1}-(J_{24,Q}+m^2I)^{-1}\right]e_1.
\]
Since \(s_j>0\),
\[
D(J_{24,Q};m)=
\sum_{j=1}^{24}u_j\left(\frac1{s_j}-\frac1{s_j+m^2}\right)>0.
\]
For any same-scheme finite remainder \(\Xi_Q<R\), set
\[
\omega_Q=\frac{3\pi(R-\Xi_Q)}{D(J_{24,Q};m)}.
\]
Then
\[
\frac{\omega_Q}{3\pi}
e_1^\top\!\left[J_{24,Q}^{-1}-(J_{24,Q}+m^2I)^{-1}\right]e_1
+\Xi_Q
=R.
\]
Distinctness makes the moment inner product nondegenerate through degree 23,
so the finite Jacobi construction has 24 sites with positive off-diagonals.
Repeated nodes instead collapse the support and do not justify a 24-site
irreducible Jacobi matrix. The construction reproduces the chosen residual exactly, but the freedom in \(s_j\), \(u_j\),
\(\omega_Q\), and \(\Xi_Q\) means it is a back-solving theorem rather than a source derivation.
\end{proof}

The source-side object needed for a final OPH endpoint theorem is therefore
\begin{equation}
\Rq(P)
=
\frac{\omega_Q(P)}{3\pi}
e_1^\top\!\left[
J_{24,Q}(P)^{-1}
-
\bigl(J_{24,Q}(P)+m_Z(P)^2I\bigr)^{-1}
\right]e_1
+\Xi_Q(P),
\label{eq:sourceJ24Q}
\end{equation}
or equivalently the source-emitted spectral density \(\rho_Q(s;P)\) plus the same-scheme finite
remainder. The required dependency graph has no path from the measured Thomson endpoint into
\(J_{24,Q}\), \(\omega_Q\), \(\Xi_Q\), or \(\rho_Q\). The available execution embeds target
constants, uses a mismatched \(P\), conflates the closure residual with the total payload, and
supplies a sampled envelope rather than a certified interval. It does not
support a physical conclusion.
A qualifying contract fixes one coordinate schema, requires a source-emitted function over the
same \(P\), certifies quadrature and interval bounds, and scores a detached artifact.
The local source-payload validator checks declared contract consistency;
it does not replay the named backend proof or establish spectral provenance.
An arbitrary nonempty certificate label cannot establish those premises;
in particular, a payload whose anchor and other components sum to an interval
disjoint from its declared total endpoint is inconsistent. Such inputs are
rejected, and a consistent but unreplayed contract is non-promoting. The separate
executable interval certificates for the two incomplete declared maps
retain their stated mathematical scope.

\section{Step 21: The Conditional Pure Source Theorem}

\begin{theorem}[OPH fine-structure endpoint, conditional pure source form]
Assume:
\begin{enumerate}[label=(\roman*)]
\item the OPH overlap-consistency branch emits the source map \(P\mapsto \AZ(P)\) by Eqs.~\eqref{eq:MU}--\eqref{eq:AZ};
\item the Ward-projected \(\U(1)_Q\) transport theorem emits a same-scheme endpoint map
\[
\ATh(P)=\AZ(P)+\Dlep(P)+\Dhad(P)+\DEW(P);
\]
\item a target-independent source law selects this source and transport map from the admitted map family;
\item a typed same-quantity bridge identifies \((P-\varphi)/\sqrt\pi\) and \(1/\ATh(P)\) as readings of one physical screen cell, with \(1/\ATh(P)\) identified as its Thomson-limit electromagnetic coupling;
\item \(G(P)=\varphi+\sqrt\pi/\ATh(P)\) is a self-map and a contraction on a nonempty closed branch interval \(I\subset\mathbb R\);
\item the interval image contains the root and the residual bound is certified.
\end{enumerate}
Then there is a unique \(P_\star\in I\) satisfying \(G(P_\star)=P_\star\), and the fine-structure constant on that branch is
\[
\alpha(0)=\frac{1}{\ATh(P_\star)}=\frac{P_\star-\varphi}{\sqrt\pi}.
\]
\end{theorem}

\begin{proof}
By assumption (v), \(I\) is nonempty and complete in the Euclidean metric,
and \(G:I\to I\) is a contraction. Banach's fixed-point theorem gives a unique \(P_\star\in I\) such that \(G(P_\star)=P_\star\). Equation~\eqref{eq:G} gives
\[
P_\star-\varphi=\frac{\sqrt\pi}{\ATh(P_\star)}.
\]
Dividing by \(\sqrt\pi\) gives
\[
\frac{P_\star-\varphi}{\sqrt\pi}=\frac{1}{\ATh(P_\star)}.
\]
The left side is \(\alpha_{\mathrm{ext}}(P_\star)\), and the right side is \(\alpha_{\mathrm{in}}(P_\star)\). Assumptions (iii)--(iv) select this map and identify their common value with the Thomson-limit electromagnetic coupling \(\alpha(0)\).
\end{proof}

Informally: no source-derived hadronic spectral endpoint map supplies the required transport. A physical fine-structure conclusion also requires target-independent map selection and the same-quantity bridge. With those premises and the interval proof, the selected full source map has one fixed point and the bridge identifies its value with the fine-structure constant.

The interval form of assumptions (v)--(vi) is direct: on any nonempty closed interval \(I\)
where interval-arithmetic evaluation
certifies \(G(I)\subseteq I\) together with a derivative bound
\(|G'|\le L<1\), the Banach fixed-point theorem gives both existence and
uniqueness of the fixed point in \(I\). The public interval contraction certificate
\cite{oph-code}
proves \(G(I)\subseteq\operatorname{int}(I)\) and \(L\le0.0724\) for both
readout modes by a direct Banach argument in mean-value (centered) form, with
\texttt{mpmath.iv} outward rounding on every elementary operation and with the
\(\SU(2)\)/\(\SU(3)\) edge-sum tails bounded by geometric majorants, so the
enclosure covers the infinite-cutoff sums. The certified unique fixed points
are \(\alpha^{-1}=136.994835177413\ldots\) for the source map (enclosure width
\(7.2\times10^{-24}\)) and \(\alpha^{-1}=137.035660136946577\ldots\) for the
gauge-width map. The certificate certifies the declared numerical map; it does
not relate either fixed point to the measured constant. Both roots lie outside
the measured comparison basin. The global at-most-one statement is
discharged on the declared analytic domain (\(\alpha^{-1}\in[100,200]\)):
\(\sup|g'|<1\) on all 256 certified pieces, both readout maps, zero
exceptional set
(domain-global uniqueness certificate in the same artifact bundle), so
each map has exactly one fixed point on the domain.
A maximal-domain extension certificate closes the exterior. An envelope
lemma built from the declared solver windows (the pixel scan window
\(\alpha_U\in[0.02,0.08]\) and the \(m_Z\) log-grid window
\(\ln(\mu_U/m_Z)\in[0,50]\)) together with the \(m_Z\)-closure identity
certifies \(1/\alpha_3(m_Z)\ge 4\) and a positive inverse-alpha readout
floor on the whole analytic domain, and an interval sweep of the full pixel
window certifies that every window-consistent fixed-point candidate of
either readout map lies inside the declared solver interval
(global uniqueness extension certificate).
Each declared map therefore has exactly one fixed point on its maximal
analytic domain, and the set of fixed points outside the declared interval
is empty; the headline enclosures above are unchanged.

A source-only theorem requires a source-derived Ward-current spectral measure, the source Jacobi
kernel, electromagnetic normalization, a same-scheme finite remainder, a
target-independent rule selecting one physical map, a typed construction that
identifies its two sides as the same quantity, a dependency graph excluding
Thomson-target leakage, a full contraction interval for the selected map, and
a prediction interval narrower than the adopted direct measurement. For the comparison standard used here, the direct rubidium recoil
result of Morel, Yao, Cladé, and Guellati-Khélifa reports
\(\alpha^{-1}=137.035999206(11)\) with 81 parts-per-trillion relative precision
\cite{morel2020}. The operational receipt is therefore
\[
\operatorname{rad}(\alpha_{\mathrm{pred}}^{-1})<1.1\times10^{-8}
\]
in inverse-alpha units for the cited direct rubidium recoil interval.

\input{../paper/tex_fragments/ENTANGLED_PAIR_ALPHA_CANDIDATE.tex}

\section{Reproducibility}

The executable fixed-point calculation is available in the public repository
\cite{oph-code}. From the code directory cited there, the main human-facing CLI is
\begin{verbatim}
python3 derive_p.py --color always
\end{verbatim}
To print only the pure source value:
\begin{verbatim}
python3 derive_p.py --no-hadron-closure
\end{verbatim}
To emit machine-readable output:
\begin{verbatim}
python3 derive_p.py --json --output runtime/report.json
\end{verbatim}
To run the raw PDG/CERN diagnostic discussed in Section~\ref{sec:raw-pdg}:
\begin{verbatim}
python3 fine_structure_fixed_point_demo.py --compare-alpha-inv 137.035999177
\end{verbatim}

Informally: the command line separates the pure source value, the CODATA/NIST comparison value, and the diagnostic values. The mathematical paper uses the same separation.

\section{Complete Source-to-Fixed-Point Chain}

\begin{longtable}{@{}L{0.18\textwidth}L{0.44\textwidth}L{0.28\textwidth}@{}}
\toprule
Step & Mathematical object & Scope \\
\midrule
\endhead
Overlap rule & Eq.~\eqref{eq:overlap} & Agreement condition; realized by native repair in the executable specification. \\
Pixel variable & \(P=a_{\mathrm{cell}}/\ell_\star^2\) & Defined. \\
Golden balance & \(\varphi=(1+\sqrt5)/2\) & Defined. \\
Boundary width & \(\alpha_{\mathrm{ext}}=(P-\varphi)/\sqrt\pi\) & Defined. \\
Inside readout & \(\alpha_{\mathrm{in}}=1/\ATh(P)\) & Defined. \\
Closure & \(P=\varphi+\sqrt\pi/\ATh(P)\) & Defined. \\
Source scale & \(M_U=E_P\e^{-2\pi}P^{1/6}\) & Defined and calculated. \\
Cell scale & \(E_{\mathrm{cell}}=E_P/\sqrt P\) & Defined and calculated. \\
Transmutation & \(v=E_{\mathrm{cell}}\exp[-2\pi/(\beta_{\mathrm{EW}}\alpha_U)]\) & Defined and calculated. \\
Gauge running & Eq.~\eqref{eq:running} & Defined and calculated; the coefficient triple is the one-loop MSSM set, a counted structural selection. \\
\(Z\)-scale self-consistency & Eq.~\eqref{eq:mZfixed} & Defined and calculated. \\
Heat-kernel closure & Eq.~\eqref{eq:pixelclosure} & Defined and calculated. \\
Electroweak anchor & \(\AZ=\alpha_{\mathrm{em}}^{-1}(m_Z^2;P)\) & Defined and calculated. \\
Charged spectrum & Eqs.~\eqref{eq:quarkmasses} and \eqref{eq:lepmasses} & Declared phenomenological transport input; not an OPH-derived mass prediction. \\
Fermion kernel & Eqs.~\eqref{eq:kernelint} and \eqref{eq:kernelclosed} & Exact one-loop kernel. \\
Screening ansatz & Eq.~\eqref{eq:Scalc} & Declared selection of the quark continuation; asserted rather than derived, with no perturbative derivation. \\
Calculated endpoint & \(A_{\mathrm{calc}}=\AZ+\Dcalc\) & Raw no-hadron diagnostic conditional on the supplied charged-spectrum continuation. \\
Bare-\(\alpha_U\) closure & \(A_{\alpha_U}^{\mathrm{fp}}\), Eq.~\eqref{eq:barealphaUaddition} & Mixed-provenance display packet; the certified gauge-width fixed point is \(137.035660136946577\ldots\), without constituting a derivation of \(\alpha\). \\
Required hadronic correction & \(\Dhadreq^{\mathrm{fp}}\), Eq.~\eqref{eq:Dempfp} & CODATA/NIST comparison residual attributed to the missing same-scheme hadronic endpoint transport. \\
Hadronic/scheme residual & \(\Rq=\ATh-A_{\mathrm{calc}}\) & Isolated endpoint contribution. \\
CODATA/NIST endpoint & \(\ATh=137.035999177(21)\) & Measured endpoint comparison value. \\
Two-\(P\) provenance & \(P_{\mathrm{root}}\) versus \(\Ppub\) & Source-branch root and measured-endpoint comparison pixel are distinct; \(\Ppub\) inherits CODATA/NIST by definition. \\
Raw PDG diagnostic & Eq.~\eqref{eq:pdgdiag} & Comparison diagnostic in a different endpoint convention. \\
Interval theorem & Banach/contraction certificate for full \(G\) & Discharged for the declared numerical map (public certificate, both readout modes); required for the full source map with the hadronic term. \\
\bottomrule
\end{longtable}

\section{Conclusion}

Conditional on target-independent map selection, the declared same-cell
bridge, and same-scheme endpoint transport, the fine-structure constant is
the coupling that makes one OPH screen cell internally consistent. The outside
reading says that the cell is displaced from the golden-ratio balance by
\((P-\varphi)/\sqrt\pi\). The inside reading sends the same \(P\) through the
declared source map, gauge closure, electroweak anchor, and Ward-projected
electromagnetic transport, producing \(1/\ATh(P)\). Neither the exact
fixed-point mathematics nor the source certificates construct that bridge or
select either certified map. The
construction fits no continuous parameter, and each declared incomplete map
has a machine-certified unique fixed point on its declared domain. The
map selector, same-cell construction, hadronic source payload, and detached
no-target-leak certificate are premises this construction does not supply, so
the residuals are diagnostics rather than a derivation of \(\alpha\).

The pure source calculation computes the root witness
\[
\alphainv_{\mathrm{root}}
=
\decnum{136.994835177412937295289429464436902857561206151035393184502},
\]
the interval-certified unique fixed point of the declared source map
(enclosure width \(7.2\times10^{-24}\), shipped interval certificate).
The CODATA/NIST measured endpoint value is
\[
\alphainv(0)=137.035999177(21),
\]
with
\[
P=1.630968209403959324879279847782648941\ldots.
\]
Combining the source root with the finite-screen unified gauge-width contribution evaluated at the
CODATA-derived comparison pixel gives the mixed-provenance no-hadron diagnostic
\[
A_{\alpha_U}^{\mathrm{fp}}=137.035959513608567790289429464437\ldots,
\]
which is not a fixed point of a single declared map. The
certified self-consistent gauge-width fixed point is
\(\alpha^{-1}=137.035660136946577\ldots\), \(2.5\times10^{-6}\) below the measured value, about
\(1.6\times10^{4}\) measurement sigma.
The residual gap is CODATA-minus-diagnostic bookkeeping, not an emitted low-energy hadronic
calculation. The physical pure-source theorem requires the Ward-projected
hadronic spectral measure, same-scheme transport, target-independent map
selection, the typed same-quantity bridge, and the interval-certificate step.
The endpoint audit holds \(P\) and the declared branch fixed. It rejects only the reading of
the source root as a complete endpoint with zero missing correction; it assigns no theory sigma,
performs no fit and does not recalibrate \(P\) \cite{ophpredictions}. The public-data audit
withholds a new prediction claim for the fine-structure constant \cite{ophpublicdata}.

\section*{Declarations}

\textbf{Funding declaration.} No funding was received for this work.

\textbf{Consent to Participate declaration.} Not applicable.

\textbf{Consent to Publish declaration.} Not applicable.

\textbf{Author Contribution declaration.} Bernhard M\"uller conceived the OPH fine-structure fixed-point formulation, prepared the mathematical derivation, wrote and checked the numerical calculation, wrote the manuscript, and prepared the reproducibility materials. Alexander Osika contributed conceptual framing, manuscript review, and approval of the final manuscript.

\textbf{Data Availability declaration.} No experimental data were generated. The manuscript cites the 2022 CODATA/NIST inverse fine-structure constant \cite{nist-value,nist-2022}, public hadronic cross-section and vacuum-polarization references \cite{pdg-xsect,dhmz,knt19}, and the OPH source materials cited below. The TeX source, OPH paper sources, and fixed-point calculation code are available in the public repository \cite{oph-code}.

\textbf{Code Availability declaration.} The fixed-point code and runtime artifacts are available in the public repository \cite{oph-code}.

\textbf{Ethics declaration.} Not applicable.

\textbf{Competing Interests declaration.} The authors declare no competing interests.

\InputIfFileExists{../paper/tex_fragments/PRIOR_WORK_FINE_STRUCTURE_CONSTANT_DERIVATION.tex}{}{}
\begin{thebibliography}{99}

\bibitem{oph-synthesis}
B.\ M\"uller, M.\ Poneder, A.\ Osika, K.\ Xue, B.\ Cassie, P.\ Nguyen, J.\ Kim, D.\ Matscheko, D.\ A.\ Friedman, J.\ Hill, W.\ T.\ Glynn, M.\ A.\ Visser, K.\ A.\ Anirudha, and B.\ de La Fourni\`ere,
\emph{Observers Who Read Their Universe into Existence: Finite Consensus, an Executable Observer, and Standard Physics}.
\url{https://github.com/FloatingPragma/observer-patch-holography/blob/r2056/flagship/from_observer_consensus_to_standard_physics.pdf}.

\bibitem{oph-smgr}
B.\ M\"uller, M.\ Poneder, A.\ Osika, K.\ Xue, P.\ Nguyen, D.\ Matscheko, and D.\ A.\ Friedman,
\emph{Standard Model Gauge Structure under Complete-Response and Observer-Transport Premises}.
GitHub source:
\url{https://github.com/FloatingPragma/observer-patch-holography/blob/r2056/paper/deriving_standard_model_gauge_structure_from_observer_overlap_consistency.tex}.
GitHub PDF:
\url{https://github.com/FloatingPragma/observer-patch-holography/blob/r2056/paper/deriving_standard_model_gauge_structure_from_observer_overlap_consistency.pdf}.

\bibitem{oph-particle}
B.\ M\"uller, M.\ Poneder, A.\ Osika, K.\ Xue, and D.\ A.\ Friedman,
\emph{Particles as Localized Topological Defects of Native Overlap Transport: Finite Theorems, a Native Decision Test and Mass Obstructions}.
GitHub source:
\url{https://github.com/FloatingPragma/observer-patch-holography/blob/r2056/paper/deriving_the_particle_zoo_from_observer_consistency.tex}.
GitHub PDF:
\url{https://github.com/FloatingPragma/observer-patch-holography/blob/r2056/paper/deriving_the_particle_zoo_from_observer_consistency.pdf}.

\bibitem{oph-code}
OPH fixed-point calculation code,
\url{https://github.com/FloatingPragma/observer-patch-holography/tree/r2056/code/P_derivation}.
Primary files:
\path{paper_math.py},
\path{derive_p.py},
\path{fine_structure_fixed_point_demo.py}.

\bibitem{oph-koide}
B.\ Mueller and B.\ de La Fourni\`ere,
\emph{The Positive-Chamber Koide Identity for Icosahedral Face Circulants}.
\url{https://github.com/FloatingPragma/observer-patch-holography/blob/r2056/extra/koide_identity_from_positive_c3_face_circulants.pdf}.

\bibitem{ophpredictions}
FloatingPragma,
\emph{Fixed observer predictions and complete-count comparison: inherited coupling comparison}, 2026.
\url{https://github.com/FloatingPragma/observer-patch-holography/blob/r2056/model/predictions/README.md}.

\bibitem{ophpublicdata}
FloatingPragma,
\emph{OPH predictions and public-data comparisons}, 2026.
\url{https://github.com/FloatingPragma/observer-patch-holography/blob/r2056/docs/research/PUBLIC_DATA_PREDICTION_AUDIT.md}.

\bibitem{ophpostdiction}
FloatingPragma,
\emph{Postdiction ledger}, 2026.
\url{https://github.com/FloatingPragma/observer-patch-holography/blob/r2056/docs/POSTDICTION_LEDGER.md}.

\bibitem{ophladder}
FloatingPragma,
\emph{The frozen-prediction ladder}, 2026.
\url{https://github.com/FloatingPragma/observer-patch-holography/blob/r2056/docs/FROZEN_PREDICTION_LADDER.md}.

\bibitem{ophpresponse}
FloatingPragma,
\emph{P-dependent electromagnetic observation and native feedback}, 2026.
\url{https://github.com/FloatingPragma/observer-patch-holography/blob/r2056/model/observer/p_response/README.md}.

\bibitem{ophentropy}
FloatingPragma,
\emph{Tangent first laws and finite entropy changes: native reference entropy and region boundaries}, 2026.
\url{https://github.com/FloatingPragma/observer-patch-holography/blob/r2056/code/geometry/ENTROPY_FIRST_LAW_AUDIT.md}.

\bibitem{ophmodel}
FloatingPragma,
\emph{OPH executable observer construction}, 2026.
\url{https://github.com/FloatingPragma/observer-patch-holography/tree/r2056/model}.

\bibitem{ophconstruction}
FloatingPragma,
\emph{OPH model construction from local repair}, 2026.
\url{https://github.com/FloatingPragma/observer-patch-holography/blob/r2056/model/CONSTRUCTION.md}.

\bibitem{jentschura-nandori}
U.\ D.\ Jentschura and I.\ Nandori,
\emph{Attempts at a determination of the fine-structure constant from first principles: A brief historical overview},
Eur.\ Phys.\ J.\ H \textbf{39}, 591--613 (2014).
\url{https://arxiv.org/abs/1411.4673}.

\bibitem{heyrovska-narayan}
R.\ Heyrovska and S.\ Narayan,
\emph{Fine-structure Constant, Anomalous Magnetic Moment, Relativity Factor and the Golden Ratio that Divides the Bohr Radius}.
\url{https://arxiv.org/abs/physics/0509207}.

\bibitem{singh}
T.\ P.\ Singh,
\emph{Quantum theory without classical time: octonions, and a theoretical derivation of the fine structure constant \(1/137\)},
Int.\ J.\ Mod.\ Phys.\ D \textbf{30}, 2142010 (2021).
\url{https://arxiv.org/abs/2110.07548}.

\bibitem{ciborowski}
J.\ Ciborowski,
\emph{Bi-Constructible pattern of weak and flavour mixing: implications for electroweak coupling constants}.
\url{https://arxiv.org/abs/2508.00030}.

\bibitem{chekanov-kjellerstrand}
S.\ V.\ Chekanov and H.\ Kjellerstrand,
\emph{Discovering the Underlying Analytic Structure Within Standard Model Constants Using Artificial Intelligence},
Particles \textbf{8}, 95 (2025).
\url{https://arxiv.org/abs/2507.00225}.

\bibitem{liu-low-yin}
Q.\ Liu, I.\ Low, and Z.\ Yin,
\emph{A Quantum Computational Determination of the Weak Mixing Angle in the Standard Model}.
\url{https://arxiv.org/abs/2509.18251}.

\bibitem{dhmz}
M.\ Davier, A.\ Hoecker, B.\ Malaescu, and Z.\ Zhang,
\emph{Reevaluation of the Hadronic Contributions to the Muon \(g-2\) and to \(\alpha(M_Z)\)},
Eur.\ Phys.\ J.\ C \textbf{71}, 1515 (2011).
\url{https://arxiv.org/abs/1010.4180}.

\bibitem{knt19}
A.\ Keshavarzi, D.\ Nomura, and T.\ Teubner,
\emph{The \(g-2\) of charged leptons, \(\alpha(M_Z^2)\) and the hyperfine splitting of muonium},
Phys.\ Rev.\ D \textbf{101}, 014029 (2020).
\url{https://arxiv.org/abs/1911.00367}.

\bibitem{morel2020}
L.\ Morel, Z.\ Yao, P.\ Cladé, and S.\ Guellati-Khélifa,
\emph{Determination of the fine-structure constant with an accuracy of 81 parts per trillion},
Nature \textbf{588}, 61--65 (2020).
\url{https://doi.org/10.1038/s41586-020-2964-7}.

\bibitem{erler-ferro}
J.\ Erler and R.\ Ferro-Hernandez,
\emph{Perturbative contributions to \(\Delta\alpha^{(5)}(M_Z^2)\)}.
\url{https://arxiv.org/abs/2308.05740}.

\bibitem{pdg-xsect}
Particle Data Group,
\emph{Data files and plots of cross-sections and related quantities in the Review of Particle Physics}.
\url{https://pdg.lbl.gov/current/xsect/}.

\bibitem{alphaqed}
F.\ Jegerlehner,
\emph{Software packages alphaQED and pQCDAdler}.
\url{https://people.physik.hu-berlin.de/~fjeger/software.html}.

\bibitem{nist-value}
NIST,
\emph{CODATA Value: inverse fine-structure constant},
2022 CODATA recommended value.
\url{https://physics.nist.gov/cgi-bin/cuu/Value?eqalphinv=}.

\bibitem{nist-2022}
P.\ J.\ Mohr, E.\ Tiesinga, D.\ B.\ Newell, and B.\ N.\ Taylor,
\emph{CODATA Recommended Values of the Fundamental Physical Constants: 2022}.
NIST publication page:
\url{https://www.nist.gov/publications/codata-recommended-values-fundamental-physical-constants-2022}.

\bibitem{nist-wall}
NIST/CODATA,
\emph{CODATA Recommended Values of the Fundamental Physical Constants: 2022},
wall chart.
\url{https://physics.nist.gov/cuu/pdf/wall_2022.pdf}.

\InputIfFileExists{../paper/tex_fragments/PRIOR_WORK_FINE_STRUCTURE_CONSTANT_DERIVATION_BIBITEMS.tex}{}{}
\bibitem{wissnergross2026pairs}
A. D. Wissner-Gross, \emph{A Parameter-Free Candidate for the Fine-Structure
Constant from Entangled Qudits}, preprint (2026).
\url{https://alexwg.org/publications/ChainTheoryAlpha.pdf}.

\end{thebibliography}

\section*{AI Assistance Disclosure}
This research project used research-grade commercial models, including
Anthropic's Fable and OpenAI's GPT-6 Astra, for research support, software
development, editing, and synthesis. The authors are responsible for the
paper's claims, methods, and final text.

\end{document}
